%%%% 原模板参考了https://www.wondercv.com/的模板
\documentclass[11pt]{article}
% disable indent globally
\setlength{\parindent}{0pt}
% some general improvements, defines the XeTeX logo
\usepackage{xltxtra}
\usepackage{bookmark}
% use hyperlink for email and url
\usepackage{hyperref}
\hypersetup{hidelinks}
\usepackage{url}
\urlstyle{tt}
\usepackage{multicol}
\usepackage{xcolor}
%%%% 统一一种颜色，偏蓝色，用于section下划线和fontawesome
\definecolor{CVBlue}{RGB}{45, 77, 152} % 从校徽上取的师大蓝
%%另一种师大蓝{RGB}{0,77,255}
%%% \widthof[]{} 用于特殊对齐是用到
\usepackage{calc}
%%%% 利用tikz来定位照片和学校Logo
\usepackage{graphicx}
\usepackage{tikz}
\usetikzlibrary{calc}
% loading fonts
\usepackage{fontspec}
\usepackage{xeCJK}
\CJKsetecglue{} %% 取消中文与数字之间间隙
%%%%% 字体需要自己下载安装，注意版权问题，这两种字体应该比较好看，英文Helvetica，中文方正兰亭黑，也是有多种版本，自己试试哪些好看。参考了https://www.wondercv.com/的模板
%%%%% windows系统好像需要先安装字体，之后下面语句就够了
%Main document font
%\setmainfont[
%  BoldFont = HelveticaNeueLTPro-Md.otf ,
%]{HelveticaNeueLTPro-Roman.otf}
%
%\setCJKmainfont[
%BoldFont=Pro_GB18030 DemiBold.otf,
%]{Pro_GB18030.otf}
%%%%% 字体需要自己下载安装，注意版权问题
%%%%% linux系统只需要字体路径就行了，如下
% % Main document font
\setmainfont[
Path = Font/,
  Extension = .otf ,
  BoldFont = HelveticaNeueLTPro-Md.otf ,
]{HelveticaNeueLTPro-Roman.otf}
\setCJKmainfont[
Path = Font/,
  Extension = .otf ,
BoldFont=ProGB18030 DemiBold.otf,
]{ProGB18030.otf}
%%%%% 定义更漂亮的“C++”，参考https://tex.stackexchange.com/questions/4302/prettiest-way-to-typeset-c-cplusplus 
%%%%% 貌似跟具体字体大小有关，需要调下参数，我测试感觉下面的比较好看
\usepackage{relsize}
\usepackage{xspace}
\protected\def\Cpp{{C\nolinebreak[4]\hspace{-.05em}\raisebox{.28ex}{\relsize{-1}++}}\xspace} 
% use fontawesome
\usepackage{fontawesome}
%\newfontfamily{\FA}{[FontAwesome.otf]}
\usepackage[
	a4paper,
	left=1.2cm,
	right=1.2cm,
	top=1.5cm,
	bottom=1cm,
	nohead
]{geometry}
\renewcommand{\baselinestretch}{1.2} %定义行间距1.2
\usepackage{titlesec}
\usepackage{enumitem}
\setlist{noitemsep} % removes spacing from items but leaves space around the whole list
%\setlist{nosep} % removes all vertical spacing within and around the list
\setlist[itemize]{topsep=0.25em, leftmargin=*}
\setlist[enumerate]{topsep=0.25em, leftmargin=*}
\titleformat{\section}         % Customise the \section command 
  {\large\bfseries\raggedright} % Make the \section headers large (\Large),
                               % small capitals (\scshape) and left aligned (\raggedright)
  {}{0em}                      % Can be used to give a prefix to all sections, like 'Section ...'
  {}                           % Can be used to insert code before the heading
  [{\color{CVBlue}\titlerule}]                 % Inserts a horizontal line after the heading
\titlespacing*{\section}{0cm}{*1.6}{*1.2}
\usepackage{siunitx}
\usepackage{amssymb}
%\xeCJKsetup{CJKspace=true}
%\xeCJKDeclareCharClass{CJK}{`0 -> `9}    % 设置 0-9 以 CJK 字体输出
%\normalspacedchars{0,1,2,3,4,5,6,7,8,9} % 0-9 的字符类被还原

% Setup fancyhdr for footer on every page
\usepackage{fancyhdr}
\pagestyle{fancy}
\fancyhf{} % clear all fields
\renewcommand{\headrulewidth}{0pt}
\renewcommand{\footrulewidth}{0pt}
\cfoot{
    \begin{tikzpicture}[remember picture, overlay]
        \node[anchor = south,fill=CVBlue,draw=none,minimum width=\paperwidth,minimum height=1.5em,align=center,font=\footnotesize,text=white] at ($(current page.south)$)
        {\centering \small 感谢您百忙中拨冗阅读我的简历！};
    \end{tikzpicture}
}

\begin{document}
%%%% Use TikZ to position the photo
\begin{tikzpicture}[remember picture, overlay]
	\node[anchor = north east] at ($(current page.north east)+(-1cm,-0.6cm)$) {\includegraphics[height=3.8cm]{img/photo.jpg}};
\end{tikzpicture}%
%%%% Use TikZ to position the university logo (shown on the first page)
\begin{tikzpicture}[remember picture, overlay]
	\node[anchor = north west] at ($(current page.north west)+(0.2cm,-0.2cm)$) {\includegraphics[height=1.75cm]{img/ustcfull.jpg}};
\end{tikzpicture}%

% The tikzpicture environment is sensitive: spaces/blank lines around comments may affect spacing.
%
\vspace*{1em}
\filright{\huge \bfseries{Zining Wang}}

\vspace*{0.4em}

\filright{\normalsize{
		\faPhone \ (+86)18221853868 \quad
		\faEnvelopeO \ \href{mailto:znwang@mail.ustc.edu.cn}{znwang@mail.ustc.edu.cn} \quad
        \faGithub \ \href{https://github.com/mogoo7zn}{https://github.com/mogoo7zn}}}
%% Prefer your .edu email address

\vspace*{0.9em}

\section{\makebox[\widthof{\faGraduationCap}][c]{\color{CVBlue}\faGraduationCap}\ Education}

\textbf{University}: University of Science and Technology of China (USTC) \\
\textbf{Program}: Undergraduate (Bachelor's Program) \\
\textbf{School}: School of Computer Science and Technology \\
\textbf{Major}: Computer Science and Technology
\newline \textbf{GPA: $\mathsf{3.6/4.3}$} \qquad \quad \textbf{Major Rank: $\mathsf{55/196}$}
\begin{itemize}[parsep=0.5ex]
	\item Relevant coursework: Compilers, Algorithms, Operating Systems, Computer Organization, etc.
    \item English: CET-4 697 \quad CET-6 641 \\
    \textbf{Able to independently read and write academic papers in English}
\end{itemize}

\vspace*{0.6em}

\section{\makebox[\widthof{\faGraduationCap}][c]{\color{CVBlue}\faTrophy}\ Honors \& Awards}
\begin{itemize}
    \item \textbf{USTC Outstanding Student Scholarship (Silver Award), 2023}
    \item \textbf{FLTRP·ETIC Cup National English Competition (Anhui Provincial Second Prize), 2023}
    \item \textbf{USTC Outstanding Student Scholarship (Silver Award), 2024}
    \item \textbf{USTC ``Yuqing Cup'' Software Development Competition (Third Prize), 2024}
    \item \textbf{National Operating System Design Contest (National Third Prize), 2025}
    \item \textbf{iGEM Competition (Gold Medal), 2025}
\end{itemize}

\vspace*{0.6em}

\section{\makebox[\widthof{\faGraduationCap}][c]{\color{CVBlue}\faUniversity}\ Research \& Project Experience}

\vspace*{0.2em}

\textbf{\large \color{CVBlue} Style Fusion for Wangjiang Tiaohua (Intangible Cultural Heritage)} \hfill 2024.01--2025.11
\begin{itemize}
    \item \textbf{Type}: Undergraduate Research Program (Machine Learning \& Software Development)
    \item \textbf{Overview}: This project focuses on the digital preservation and innovation of Wangjiang Tiaohua, a Chinese intangible cultural heritage.
    By fusing the aesthetics of traditional craftsmanship with user-customized images, we generate automated knitting/embroidery patterns and build a new digital creation workflow.
    The approach reduces the design workload and time cost for inheritors while enabling batch generation of pattern proposals that balance traditional charm with modern aesthetics for potential commercial applications.
    The work was conducted collaboratively; I primarily implemented style fusion and image-processing components and supported teammates in developing the knitting-process software system.

    The system has two key innovations: (1) intelligent image parsing that converts raw materials into operable knitting patterns; (2) a process-decomposition algorithm that generates step-by-step visual instructions to precisely guide production.
    The project was carried out under the guidance of university faculty and heritage inheritors to ensure both academic rigor and cultural authenticity.

    \item[\textcolor{CVBlue}{\faStar}] \textbf{\textcolor{CVBlue}{Highlights}}: \textbf{Patent application in progress}; \quad Hands-on experience in fine-tuning diffusion models and dataset annotation
    \item \textbf{Tech Stack}: Python, PyTorch, Computer Vision, Diffusion Models; C\#, WPF
\end{itemize}

\par\noindent\textcolor{black!20}{\rule{\linewidth}{0.5pt}}

\vspace*{0.1em}

\textbf{\large \color{CVBlue} National College Student Operating System Competition} \hfill 2025.04--2025.08
\begin{itemize}
    \item \textbf{Type}: Academic Competition
    \item \textbf{Overview}: In the preliminary round, we deployed a large language model on a DAYU200 development board under limited CPU and memory resources, built an on-device LLM application, and ran benchmark tests.
    In the final round, we developed an application on Huawei devices based on HarmonyOS, leveraged device GPU compute to improve inference speed, and explored early integration of multimodal-LLM capabilities with the operating system.
    We achieved approximately 15 token/s for on-device multimodal inference deployment. I served as the team leader.
    \item[\textcolor{CVBlue}{\faStar}]\textbf{\textcolor{CVBlue}{Core Engineering}}: \textbf{Embedded development and deployment of a multimodal LLM application}
    \item \textbf{Tech Stack}: C++, ArkTS, Llama, MNN
\end{itemize}

\par\noindent\textcolor{black!20}{\rule{\linewidth}{0.5pt}}

\vspace*{0.1em}

\textbf{\large \color{CVBlue} Huawei HarmonyOS Elite Program Internship} \hfill 2025.06--2025.08
\begin{itemize}
    \item \textbf{Type}: Industry Internship

    \item \textbf{Overview}: Selected through on-campus screening to join Huawei's HarmonyOS Elite Program, where I studied OpenHarmony and HarmonyOS technologies and completed an internship with an internal team.
    I developed an automated parameter adaptation tool for a camera subsystem team to improve parameter-tuning efficiency and reduce adaptation and version-iteration costs.
    \item[\textcolor{CVBlue}{\faStar}] \textbf{\textcolor{CVBlue}{Key Takeaway}}: Gained understanding of industry workflows and learned that, regardless of role, the most critical practical ability is to quickly locate and solve problems \emph{and} communicate the solution clearly.
    \item \textbf{Tech Stack}: C++, Python
\end{itemize}

\par\noindent\textcolor{black!20}{\rule{\linewidth}{0.5pt}}

\vspace*{0.1em}

\textbf{\large \color{CVBlue} Kaggle -- ConnectX} \hfill 2025.10--2025.12
\begin{itemize}
    \item \textbf{Type}: Course Project

    \item \textbf{Overview}: Built agents for gravity-based Connect Four using reinforcement learning methods such as DQN and AlphaZero, and competed in Kaggle ConnectX.
    Best historical rank: 19/241.
    \item[\textcolor{CVBlue}{\faStar}] \textbf{\textcolor{CVBlue}{Outcome}}: Learned foundational RL methods and math principles; gained practical experience developing RL agents
    \item \textbf{Tech Stack}: PyTorch, DQN, AlphaZero, RL
\end{itemize}

\par\noindent\textcolor{black!20}{\rule{\linewidth}{0.5pt}}

\vspace*{0.1em}

\textbf{\large \color{CVBlue} TimeFlow: Intelligent Time Management App} \hfill 2023.10--2024.11
\begin{itemize}
    \item \textbf{Type}: Android App Development

    \item \textbf{Overview}: Developed a time-optimization management app for students and faculty.
    Based on user habit analysis, the system combines personal routines with a task-priority matrix to automatically generate more scientific schedules and iteratively refine planning strategies, providing a solution for maximizing time value in educational settings.
    \item[\textcolor{CVBlue}{\faStar}] \textbf{\textcolor{CVBlue}{Outcome}}: Learned the Android mobile development workflow and frameworks; gained an end-to-end view from user needs to implementation, feedback, and iterative improvement
    \item \textbf{Tech Stack}: Kotlin, Java, MySQL
\end{itemize}

\par\noindent\textcolor{black!20}{\rule{\linewidth}{0.5pt}}

\vspace*{0.1em}

\textbf{\large \color{CVBlue} iGEM (International Genetically Engineered Machine) Competition} \hfill 2025.02--2025.10
\begin{itemize}
    \item \textbf{Type}: Academic Competition
    \item \textbf{Overview}: Co-built an on-campus experimental server platform with web-team members and designed the competition website (UI/UX and functionality).
    Learned web front-end and back-end development, and more importantly gained initial experience leading a team: planning task allocation, coordinating with other subteams, and handling conflicts and urgent issues.
    Served as the web-team lead.
    \item[\textcolor{CVBlue}{\faStar}] \textbf{\textcolor{CVBlue}{Outcome}}: Web development skills; \ \textbf{improved team communication and management—learning to collaborate, negotiate, and encourage rather than working solo}
    \item \textbf{Tech Stack}: JavaScript, React, Next.js, Haskell
\end{itemize}

\vspace*{0.7em}

\section{\makebox[\widthof{\faGraduationCap}][c]{\color{CVBlue}\faUsers}\ Leadership \& Service}
\begin{itemize}
	\item 2025 \quad Youth League Branch Secretary, University Flag Guard Team
	\item 2023--2025 \quad Class Life Committee Member
\end{itemize}

\vspace*{0.7em}

\section{\makebox[\widthof{\faGraduationCap}][c]{\color{CVBlue}\faWrench}\ Technical Skills}

\begin{itemize}
    \item \textbf{Platforms}: Embedded boards (OrangePi AI Pro, DAYU200); OpenHarmony/HarmonyOS; Android; WPF 
    \item \textbf{Languages}: C++, Java, C\#, Kotlin, Python, ArkTS
    \item \textbf{Frameworks \& Libraries}: PyTorch, Next.js
    \item \textbf{Tools}: Git, Docker, Anaconda, \LaTeX{}, Android Studio
\end{itemize}


%%%% If the resume spans multiple pages, insert \newpage or \clearpage manually where needed.


\end{document}
